\documentclass{minmal}
\usepackage{array}
\renewcommand{\arraystretch}{1.5}
\begin{document}
\begin{center}
    \textbf{Prøve i mekaniske separasjonsmetoder:}
\end{center}
\noindent
\textbf{Oppg. 1:}\\\\
\noindent
a) Hva menes en homogen og en heterogen blanding? Gi et eksempel for hver av dem.\nl[1]
b) Hva menes med en suspensjon? Gi et eksempel på en suspensjon. \nl[1]
c) Hva vil skje med partikler i vann som har større tetthet enn vann, og hva vil skje med partikler i vann som har tetthet mindre enn vann?\nl[1]
d) Hvilke to emulsjoner har vi? 
e) Hva menes med mettet, umettet og overmettet løsning. Bruk $NaCl$ som eksempel.  Ved $25^0C$ for du maksimum løst 360g/L.  \nl[2]

\noindent \textbf{Oppg. 2:}\\\\
a) Figuren under viser en standard rørebeholder:\\

\begin{tabular}{m{8cm}>{\raggedright\arraybackslash}m{6cm}}

\includegraphics[width=8cm]{rorebeholder.png} &
1) Hva brukes inspeksjonsluken til? \newline \newline
2) Hva er oppgaven til temperaturføleren? \newline \newline
3) Hva er oppgaven til elektromotoren? \\

\end{tabular}\newpage\noindent
b) Forklar kort hva som menes med begrepet viskositet for en væske. Gi et eksempel på en væske med høy viskositet og en med lav viskositet.\nl[1]
c) Det finnes mange typer rørere:
\begin{figure}[h]
    \centering
    \includegraphics[width=0.9\linewidth]{rørere.png}
    \caption{Ulike rørere}
    \label{fig:placeholder}
\end{figure}
En propellrører er hurtiggående, en spiral Wendel er saktegående. hvilken av de to rørerne vil du bruke i en sirup-tykk løsning? Begrunn svaret.\nl[1]
\textbf{Oppg. 3:}\nl[1]
a) Forklar forskjellen på sortering og klassering i kjemisk teknologi.\nl[1]
b) Forklar virkemåten til en hydrosyklon:
\begin{figure}[h]
    \centering
    \includegraphics[width=0.65\linewidth]{hydrosyklon.png}
    \caption{En hydrosyklon}
    \label{fig:placeholder}
\end{figure}\newpage
c) Forklar virkemåten til magnetsortereren:\\
\begin{figure}[h]
    \centering
    \includegraphics[width=0.7\linewidth]{magnetseparator.png}
    \caption{magnetsorterer}
    \label{fig:placeholder}
\end{figure}\\\
d) Forklar virkemåten til et rundt klaringsbasseng:
\begin{figure}[h]
    \centering
    \includegraphics[width=0.84\linewidth]{rundt_klaringsbasseng.png}
    \caption{Et rundt klaringsbasseng}
    \label{fig:placeholder}
\end{figure}\newpage
e) Forklar virkemåten til flotasjonsfortykkeren:
\begin{figure}[h]
    \centering
    \includegraphics[width=0.9\linewidth]{flotasjonsfortykker.png}
    \caption{Flotasjonsfortykker}
    \label{fig:placeholder}
\end{figure}\nl[1]
f) For rensing av gass bruker man et støvkammer:
\begin{figure}[h]
    \centering
    \includegraphics[width=0.8\linewidth]{stovkammer.png}
    \caption{Støvkammer}
    \label{fig:placeholder}
\end{figure}\nl[1]
Forklar virkemåten til støvkammeret.
\end{document}